\documentclass[11pt]{article}
\usepackage[margin=1in]{geometry}
\usepackage{amsmath,amssymb}

\title{Course Notes for MATH 375: Probability Theory}
\author{Mariana Restrepo}
\date{Fall 2026}

\begin{document}
\maketitle

\begin{center}
\large Lecture 1
\end{center}

\tableofcontents

\section{Lecture 1: Sets and Counting}

\subsection{Union and intersection}

The \textbf{union} $A\cup B$ contains every element that belongs to
$A$, to $B$, or to both:
\[
A\cup B=\{x:x\in A\text{ or }x\in B\}.
\]

The \textbf{intersection} $A\cap B$ contains the elements shared by
both sets:
\[
A\cap B=\{x:x\in A\text{ and }x\in B\}.
\]

For example, if
\[
A=\{1,2,5,7\}, \qquad B=\{2,3,6,8\},
\]
then
\[
A\cup B=\{1,2,3,5,6,7,8\},
\qquad
A\cap B=\{2\}.
\]

\subsection{Special sets and set properties}

The \textbf{universal set} $U$ contains all elements under discussion.
The \textbf{empty set} $\varnothing$ contains no elements.

Some useful identities are
\[
\begin{aligned}
A\cup B &= B\cup A, & A\cap B &= B\cap A,\\
A\cup\varnothing &= A, & A\cap\varnothing &= \varnothing,\\
A\cup U &= U, & A\cap U &= A.
\end{aligned}
\]

The distributive laws are
\[
(A\cup B)\cap C=(A\cap C)\cup(B\cap C)
\]
and
\[
(A\cap B)\cup C=(A\cup C)\cap(B\cup C).
\]

\subsection{Complements and De Morgan's laws}

The \textbf{complement} $A^c$ contains the elements of $U$ that are
not in $A$:
\[
A^c=\{x\in U:x\notin A\}.
\]

De Morgan's laws are
\[
(A\cup B)^c=A^c\cap B^c,
\qquad
(A\cap B)^c=A^c\cup B^c.
\]

\subsection{Cardinality and inclusion--exclusion}

For a finite set $A$, its \textbf{cardinality} $|A|$ is its number
of elements. To count the elements in a union, use
\[
|A\cup B|=|A|+|B|-|A\cap B|.
\]
We subtract $|A\cap B|$ because adding $|A|$ and $|B|$ counts
the shared elements twice.

For three sets,
\[
\begin{aligned}
|A\cup B\cup C|
={}&|A|+|B|+|C|\\
&-|A\cap B|-|A\cap C|-|B\cap C|\\
&+|A\cap B\cap C|.
\end{aligned}
\]

\subsection{The multiplication principle}

If a task has several successive choices, multiply the number of
options at each step.

\paragraph{Bit strings.}
Each position in a bit string can be either $0$ or $1$. Therefore,
the number of bit strings of length $n$ is
\[
\underbrace{2\cdot2\cdots2}_{n\text{ positions}}=2^n.
\]

\paragraph{License plates.}
Suppose a plate has four digits followed by three letters, and
repetition is allowed. There are
\[
10^4\cdot26^3
\]
possible plates.

If the four digits must be distinct and the three letters must be
distinct, the count becomes
\[
(10\cdot9\cdot8\cdot7)(26\cdot25\cdot24)
=78{,}624{,}000.
\]

\paragraph{Area codes.}
If an area code has three digits and its first digit cannot be $0$
or $1$, while the other digits have no restrictions, there are
\[
8\cdot10\cdot10=800
\]
possible area codes.

\subsection{Choosing positions: combinations}

How many bit strings of length $n$ have exactly one $0$?
Choose its position. There are $n$ choices.

How many have exactly two $0$s? Choose two of the $n$ positions:
\[
\binom{n}{2}
=\frac{n!}{2!(n-2)!}
=\frac{n(n-1)}{2}.
\]

For instance, with $n=5$,
\[
\binom{5}{2}=10.
\]

\textbf{Decision cue:} When identical special objects must occupy
certain positions, choose the positions. You do not need to arrange
the remaining identical objects separately.

\subsection{Permutations}

A \textbf{permutation} of $n$ distinct objects is an ordering of
all $n$ objects. The number of such orderings is
\[
n(n-1)(n-2)\cdots2\cdot1=n!.
\]

\textbf{Decision cue:} If changing the order creates a different
outcome, think about permutations. If you only need to select which
positions or objects are involved, think about combinations.

\end{document}